% ============================================================
%  §3.1 Συνδυαστικό μοντέλο πρόσφυσης
% ============================================================

\label{subsec:combined_slip}

\subsubsection{Μεγέθη ολίσθησης}

Πριν από την παρουσίαση του μοντέλου, απαιτείται ο αυστηρός ορισμός των βασικών μεγεθών
ολίσθησης. Η έννοια της «ολίσθησης» εκφράζει την απόκλιση μεταξύ της πραγματικής κίνησης
του ελαστικού και της «ιδανικής» κίνησης χωρίς ολίσθηση.

\paragraph{Γωνία ολίσθησης $\alpha$ (Slip Angle).}
Ορίζεται ως η γωνία μεταξύ του επιπέδου του τροχού (κατεύθυνση κύλισης) και του
πραγματικού διανύσματος ταχύτητας του κέντρου τροχού:

\begin{equation}
    \alpha = -\arctan\!\left(\frac{v_{y}^{\text{ελαστικού}}}{v_{x}^{\text{ελαστικού}}}\right)
    \label{eq:slip_angle}
\end{equation}

όπου $v_{x}^{\text{ελαστικού}}$, $v_{y}^{\text{ελαστικού}}$ οι συνιστώσες ταχύτητας στο επίπεδο
του πέλματος, με διεύθυνση $x$ τη διεύθυνση κύλισης και $y$ κάθετη σε αυτή. Σε καθαρή κύλιση 
$\alpha = 0$. Μη μηδενική $\alpha$ σημαίνει ότι ο τροχός
«γλιστράει» πλευρικά -- έχει εγκάρσια σχετική ταχύτητα με το οδόστρομα -- 
και αναπτύσσει πλευρική δύναμη $F_y$ (cornering force).

% ----- TikZ: slip angle -----
\begin{figure}[htbp]
\centering
\begin{tikzpicture}[scale=1.25,
    wheel/.style={rectangle,draw=black,fill=gray!30,
                  minimum width=0.32cm,minimum height=1.1cm,rounded corners=2pt},
    vec/.style={->,thick}
]
  % wheel
  \node[wheel] (W) at (0,0) {};
  % heading (rolling direction)
  \draw[vec,black!60] (0,-0.9)--(0,1.3)
    node[above,font=\small]{κατεύθυνση τροχού};
  % actual velocity
  \draw[vec,blue!80!black,very thick] (0,0)--(0.65,1.05)
    node[right,font=\small,blue!80!black]{$\mathbf{v}_{\rm wheel}$};
  % alpha arc
  \draw[->,thin,red] (0,0.52) arc[start angle=90,end angle=58,radius=0.52];
  \node[red,font=\small] at (0.3,0.72) {$\alpha$};
  % components
  \draw[dashed,gray!70] (0,1.05)--(0.65,1.05);
  \draw[vec,gray!60] (0,0)--(0.65,0) node[below,font=\scriptsize,gray!60]{$v_y$};
  \draw[vec,gray!60] (0.65,0)--(0.65,1.05) node[right,font=\scriptsize,gray!60]{$v_x$};
\end{tikzpicture}
\caption{Ορισμός γωνίας ολίσθησης $\alpha$. Η διαφορά μεταξύ της πραγματικής ταχύτητας $\mathbf{v}$ και της κατεύθυνσης κύλισης δίνει την $\alpha$, η οποία αντιστοιχεί σε πλευρική δύναμη $F_y$.}
\label{fig:slip_angle}
\end{figure}
% ----- τέλος TikZ -----


\paragraph{Λόγος ολίσθησης $\kappa$ (Slip Ratio).}
Εκφράζει τη σχετική διαφορά μεταξύ της περιφερειακής ταχύτητας τροχού και της ταχύτητας
οχήματος:
\begin{equation}
    \kappa = \frac{\omega \cdot r_{\mathrm{eff}} - v_x}{v_x}
    \label{eq:slip_ratio}
\end{equation}

Για $\kappa = 0$ υπάρχει ιδανική κύλιση (χωρίς ολίσθηση). Θετικές τιμές ($\kappa > 0$) αντιστοιχούν σε
κινητήρια ροπή ενώ αρνητικές ($\kappa < 0$) σε φρενάρισμα.

\subsubsection{Μοντέλο Pacejka -- Combined Slip}

Όπως αναφέρθηκε και στη θεωρία, τα περισσότερα μαθηματικά μοντέλα ελαστικών 
παίρνουν τη μορφή των (\cref{eq:longitudinal_tyre_force, eq:lateral_tyre_force}). 
Η συγκεκριμένη επιλογή εξαρτάται απο πολλούς παράγοντες, π.χ αν η συμπεριφορά του ελαστικού
εμφανίζει έντονες μη-γραμμικότητες ή απο τα φαινόμενα που θέλουμε να αντικατοπτρίσουμε στο μοντέλο.
Στην παρούσα διπλωματική εργασία επιλέχθηκε ένα σχετικά απλό μοντέλο combined slip 
των \textit{Tremblett & Limebeer} \cite{Tremlett2016} αφενός διότι περιγράφει με ικανοποιητική
μορφή τα φαινόμενα που θα εξεταστούν, και αφετέρου, διότι επιτρέπει την άμεση σύγκριση των αποτελεσμάτων.

Ωστόσο, η προσέγγιση είναι το μοντέλο των ελαστικών να είναι ένα τμήμα αποσπασμένο
από το γενικότερο μοντέλο του οχήματος (modular) έτσι ώστε σε διαφορετική εφαρμογή, όχημα ή ελαστικό
να μπορεί να αντικατασταθεί απο μοντέλο που ταιριάζει καλύτερα στις ανάγκες του ή στα 
πειραματικά δεδομένα.

\paragraph{Φαινόμενο load sensitivity.}
Ο συντελεστής τριβής $\mu$ δεν είναι σταθερός αλλά μειώνεται ελαφρά καθώς αυξάνεται το
κατακόρυφο φορτίο. Το φαινόμενο αυτό -- γνωστό ως \emph{load sensitivity} -- μοντελοποιείται
με γραμμική παρεμβολή:

\begin{equation}
    \mu(F_z) = \mu_1 + \frac{\mu_2 - \mu_1}{F_{z2} - F_{z1}}\,(F_z - F_{z1}),
    \quad F_{z1} = 2000\,\mathrm{N},\; F_{z2} = 6000\,\mathrm{N}
    \label{eq:mu_load}
\end{equation}

Η φυσική βάση αυτού του φαινομένου έγκειται στη μη-γραμμική ελαστικότητα του ελαστομερούς
υπό πίεση: σε μεγάλα φορτία ο μηχανισμός προσκόλλησης φτάνει σε κορεσμό. Πρακτικά, αυτό σημαίνει
ότι κατά τη μεταφορά φορτίου σε μια στροφή ο επιφορτισμένος τροχός κερδίζει λιγότερη πρόσφυση
από ό,τι χάνει ο αποφορτισμένος. Αυτός ο μηχανισμός είναι και στο μεγαλύτερο βαθμό ο λόγος για τον
οποίο επιθυμείται χαμηλό κέντρο βάρους σε αγωνιστικές εφαρμογές (σχέσεις \cref{eq:Fz_long}, \cref{eq:Fz_lat})
-- ελαχιστοποιώντας το ύψος του κέντρου βάρους ελαχιστοποιούμε και το ποσό πρόσφυσης που "χάνεται".

\section{ΕΔΩ ΕΧΩ ΜΕΙΝΕΙ}

\paragraph{Κανονικοποιημένη ολίσθηση και συνδυασμένη ολίσθηση.}
Ορίζεται $\kappa_n \in [-1,1]$ με $\kappa = \kappa_n \cdot \kappa_{\max}$ ώστε να
αποφευχθούν αριθμητικά προβλήματα. Ο συνδυαστικός δείκτης ολίσθησης:
\begin{equation}
    \sigma = \sqrt{\kappa_n^2 + \tan^2\!\alpha}
    \label{eq:combined_slip}
\end{equation}
χρησιμοποιείται για τον υπολογισμό της συνολικής δύναμης τριβής, η οποία στη συνέχεια
κατανέμεται σε $F_x$ και $F_y$ ανάλογα με τις συνιστώσες $\kappa_n$ και $\tan\alpha$.

\paragraph{Εξάρτηση με θερμοκρασία.}
\begin{equation}
    \mu_{\rm eff}(F_z, T) = \mu(F_z) \cdot g(T)
    \label{eq:mu_temp}
\end{equation}
όπου $g(T)$ η συνάρτηση scaling πρόσφυσης θερμοκρασίας: αυξάνεται στο thermal window και
μειώνεται εκτός αυτού. Η επιλογή της μορφής του $g(T)$ (τετραγωνικό, Gaussian, κλπ.)
επηρεάζει πόσο απότομα εμφανίζεται το cliff πρόσφυσης.

% TODO: Ορίσε ρητά τη μορφή του g(T) που χρησιμοποιείς στον κώδικα
% (getGripScalingFactorFromTemperature). Αξίζει να δείξεις το γράφημα g(T) vs T.